\documentclass[12pt,twoside]{article}
%\usepackage[table]{xcolor} % important to avoid options clash.
%\input{02465shared_preamble}
%\usepackage{cleveref}
\usepackage{url}
\usepackage{graphics}
\usepackage{multicol}
\usepackage{rotate}
\usepackage{rotating}
\usepackage{booktabs}
\usepackage{hyperref}
\usepackage{pifont}
\usepackage{latexsym}
\usepackage[english]{babel}
\usepackage{epstopdf}
\usepackage{etoolbox}
\usepackage{amsmath}
\usepackage{amssymb}
\usepackage{multirow,epstopdf}
\usepackage{fancyhdr}
\usepackage{booktabs}
\usepackage{xcolor}
\newcommand\redt[1]{ {\textcolor[rgb]{0.60, 0.00, 0.00}{\textbf{ #1} } } }


\newcommand{\m}[1]{\boldsymbol{ #1}}
\newcommand{\EE}{\mathbb{E}}
\DeclareMathOperator*{\argmax}{arg\,max}
\DeclareMathOperator*{\argmin}{arg\,min}
\newcommand{\yoursolution}{ \redt{(your solution here) } } 



\title{ Report 3 hand-in }
\date{ \today }
\author{Alice (\texttt{s000001})\and  Bob (\texttt{s000002})\and Clara (\texttt{s000003}) } 

\begin{document}
\maketitle

\begin{table}[ht!]
\caption{Attribution table. Feel free to add/remove rows and columns}
\begin{tabular}{llll}
\toprule
                                                                           & Alice   & Bob    & Clara   \\
\midrule
 1: Optimal policy                                                         & 0-100\%  & 0-100\% & 0-100\%  \\
 2: Simulating a finite approximation of the optimal action-value function & 0-100\%  & 0-100\% & 0-100\%  \\
 3: Analytically computing the optimal action-value function               & 0-100\%  & 0-100\% & 0-100\%  \\
 4: Extend solution to all states and actions                              & 0-100\%  & 0-100\% & 0-100\%  \\
 5: Modified reward function                                               & 0-100\%  & 0-100\% & 0-100\%  \\
 6: Long bridge                                                            & 0-100\%  & 0-100\% & 0-100\%  \\
\bottomrule
\end{tabular}
\end{table}

%\paragraph{Statement about collaboration:}
%Please edit this section to reflect how you have used external resources. The following statement will in most cases suffice: 
%\emph{The code in the irls/project1 directory is entirely}

%\paragraph{Main report:}
Headings have been inserted in the document for readability. You only have to edit the part which says \yoursolution. 

\section{Jar-Jar at the battle of Naboo (\texttt{jarjar.py})}
\subsubsection*{{\color{red}Problem 3:  Analytically computing the optimal action-value function}}
	
		Using that ... we obtain 
		\begin{align}
			Q^*(0,1) & = \cdots \\
			Q^*(1,-1) & = \cdots
		\end{align}
		therefore...
	
\subsubsection*{{\color{red}Problem 5:  Modified reward function}}
	
Using ... we get that ...:
		\begin{align}
			v_\pi(s) = \dots.
		\end{align}
	
\section{Optimal escape}
\subsubsection*{{\color{red}Problem 6:  Long bridge}}
	
			Because ... we see that if $\cdots > \cdots$ one should go left, if  $\cdots < \cdots$ one should go right,
	and when $\cdots = \cdots$ it makes no difference.
		
\end{document}